\documentclass[11pt]{article}

\usepackage[utf8]{inputenc}
\usepackage[T1]{fontenc}
\usepackage{lmodern}
\usepackage{geometry}
\usepackage{amsmath,amssymb,amsfonts,amsthm,mathtools}
\usepackage{bm}
\usepackage{enumitem}
\usepackage{microtype}
\usepackage{hyperref}
\usepackage{titlesec}
\usepackage{booktabs}
\usepackage{array}
\usepackage{setspace}
\usepackage{mathrsfs}
\usepackage{physics}

\geometry{margin=1in}
\hypersetup{colorlinks=true, linkcolor=black, urlcolor=blue, citecolor=black}
\setlist[enumerate]{topsep=0.4em,itemsep=0.35em}
\setlist[itemize]{topsep=0.3em,itemsep=0.25em}
\titleformat{\section}{\large\bfseries}{\thesection.}{0.5em}{}
\titleformat{\subsection}{\normalsize\bfseries}{\thesubsection.}{0.5em}{}
\setstretch{1.06}

\newcommand{\Aether}{\AE{}ther}
\newcommand{\Aflow}{\AE{}ther-flow}
\newcommand{\TheoryName}{\Aflow{} Interpretation of Relativity}
\newcommand{\TheoryProgram}{\Aether{} / \Aflow{} framework}

\newcommand{\CompletedMetric}{g^{\sharp}}
\newcommand{\CompletionCore}{\mathfrak C^{\mathrm{zr,aff}}}
\newcommand{\TraceCore}{\mathfrak C^{\mathrm{tr}}}
\newcommand{\AffineNucleus}{\mathfrak N^{\mathrm{aff}}}

\newcommand{\Car}{\mathrm{car}}
\newcommand{\Cur}{\mathrm{cur}}
\newcommand{\Hom}{\mathrm{hom}}

\newcommand{\ForSpace}[1]{\mathcal I_{#1}^{\mathrm{for}}}
\newcommand{\PrimShell}{\mathcal S_{\mathrm{prim}}}
\newcommand{\Reduction}[1]{\mathcal R_{#1}}
\newcommand{\CurrentCarrier}[1]{\mathfrak C_{#1}^{\mathrm{cur}}}
\newcommand{\EqCarrier}[1]{\mathfrak C_{#1}^{\mathrm{eq}}}
\newcommand{\MetricOut}{\mathscr G_{\mathrm{out}}}
\newcommand{\MatterOut}{\mathscr T_{\mathrm{out}}}

\newcommand{\VisibleAffine}[1]{\widehat X_{#1}}
\newcommand{\DataSpace}[1]{\mathcal Y_{#1}^{\mathrm{obs}}}
\newcommand{\AmbientTarget}[1]{E_{#1}^{\mathrm{amb}}}

\newcommand{\BalSpace}[1]{\mathcal Z_{#1}}
\newcommand{\BalBody}[1]{\mathcal B_{#1}^{\mathrm{bal}}}
\newcommand{\BalData}[1]{\Lambda_{#1}^{\mathrm{dat}}}
\newcommand{\BalShell}[1]{\Lambda_{#1}^{\mathrm{sh}}}
\newcommand{\ZeroHiddenSlice}[1]{\mathcal Z_{#1}^{0}}
\newcommand{\ZeroHiddenVec}[1]{V_{#1}^{0}}
\newcommand{\Kernel}[1]{K_{#1}^{0}}
\newcommand{\BalOp}[1]{\mathcal K_{#1}^{\mathrm{bal}}}
\newcommand{\BalSym}[1]{\mathbb B_{#1}}
\newcommand{\BalTime}[1]{\vartheta_{#1}^{\mathrm{bal}}}
\newcommand{\BalEmbed}[1]{\zeta_{#1}^{\mathrm{bal}}}
\newcommand{\ReferencePoint}[1]{z_{#1}^{\ast}}
\newcommand{\StateComp}[1]{P_{#1}^{\mathrm{co}}}

\newcommand{\GapSpace}[1]{H_{#1}}
\newcommand{\GapBody}[1]{D_{#1}}
\newcommand{\HiddenOperator}[1]{K_{#1}^{\mathrm{hid}}}

\newcommand{\CoreState}[1]{\mathcal C_{#1}}
\newcommand{\TraceBody}[1]{P_{#1}^{\mathrm{tr}}}
\newcommand{\TraceCoord}[1]{\pi_{#1}^{\mathrm{tr}}}
\newcommand{\TraceDataMap}[1]{\Omega_{#1}^{\mathrm{tr,dat}}}
\newcommand{\TraceShellMap}[1]{\Omega_{#1}^{\mathrm{tr,sh}}}
\newcommand{\TraceSym}[1]{\mathbb K_{#1}^{\mathrm{tr}}}
\newcommand{\TraceTime}[1]{\tau_{#1}^{\mathrm{tr}}}
\newcommand{\TraceEmbed}[1]{\zeta_{#1}}

\newcommand{\TraceShellLin}[1]{L_{#1}^{\mathrm{tr,sh}}}
\newcommand{\ShellKer}[1]{V_{#1}^{\mathrm{ker}}}
\newcommand{\ShellNormal}[1]{E_{#1}^{\mathrm{sh}}}
\newcommand{\ShellOffset}[1]{\lambda_{#1}^{\mathrm{sh}}}
\newcommand{\DirProj}[1]{\Pi_{#1}^{\parallel}}
\newcommand{\CoordSpace}[1]{\mathcal K_{#1}}
\newcommand{\CoordBody}[1]{P_{#1}}
\newcommand{\CoordIso}[1]{T_{#1}}
\newcommand{\AffineData}[1]{\widehat D_{#1}}
\newcommand{\ShellAffine}[1]{\mathcal A_{#1}^{\mathrm{sh}}}
\newcommand{\ShellEmbed}[1]{\chi_{#1}}
\newcommand{\SymGroup}[1]{\mathbb G_{#1}}
\newcommand{\TimeRev}[1]{\vartheta_{#1}}

\newtheorem{definition}{Definition}
\newtheorem{proposition}{Proposition}
\newtheorem{remark}{Remark}
\newtheorem{corollary}{Corollary}
\newtheorem{theorem}{Theorem}

\title{\textbf{The \TheoryName{}}\\[0.4em]
\large Primitive-Reservoir Observer-Reduction\\
\large Minimal Completion-State Shell-Trace Core\\
\large Necessity / Uniqueness from Completion-Balance Ambient Dynamics}
\author{}
\date{}

\begin{document}

\maketitle

\begin{abstract}
The current active theorem on the primitive-reservoir observer-reduction line has already sharpened the live internal burden below the completion-core frontier: the full minimal zero-response affine completion-core presentation \(\CompletionCore\) carries no extra benchmark-facing freedom beyond a smaller completion-state shell-trace core \(\TraceCore\). The remaining honest benchmark-facing move is therefore no longer to restate the larger completion-core graph package, but to derive or uniquely force that smaller trace core from still more primitive explicit substrate structure, or else collapse it further.

The present note takes the preferred first route. For each sector it introduces the same completion-balance ambient dynamics package already preserved on the active line: a compact convex ambient admissible body, affine observer-data and shell-response readouts, an invariant affine zero-hidden slice, a positive ambient balance operator whose kernel is exactly the zero-hidden translation space, exact splitting of the ambient body over zero-hidden and hidden directions, symmetry and time-reversal actions preserving the whole package, an exact zero-response shell realization, fixed-data shell solvability on the zero-hidden slice, and a coercive balance inequality on data-invisible zero-hidden directions modulo the inert completion kernel. From that package it derives directly the current live datum \(\TraceCore\): the faithful completion-state quotient of the zero-hidden slice, its compact convex state body, the descended affine observer-data and shell-response trace maps, the hidden body with positive gap operator, the induced symmetry and time-reversal structure, the exact zero-response trace realization, fixed-data solvability, and trace-kernel coercive control.

This is a direct forcing theorem for the current live burden rather than another restatement of a preserved deeper relay. The shell-coordinate construction yielding the current exact-shell affine nucleus \(\AffineNucleus\) already factors through the ambient-induced trace core, so the trace core is no longer an independently unexplained stopping point on the active line. The benchmark package remains fixed throughout: the operative metric is still exactly \(\CompletedMetric\), the ordinary matter route is unchanged, there is no second operative metric, and the low-energy matter law is not enlarged.
\end{abstract}

\section{Introduction}

The active derivational program of \TheoryName{} has again tightened what counts as an honest next step. The direct exact-shell-affine-nucleus forcing theorem already removed \(\AffineNucleus\) as the live unexplained stopping point by deriving it canonically from the explicit completion-core package \(\CompletionCore\). The later trace-core collapse theorem then sharpened that burden once more by showing that the full affine completion-plane and completion-body presentation carries no extra benchmark-facing freedom beyond the smaller completion-state shell-trace core \(\TraceCore\).

That sharper routing result changes what a deeper-origin theorem must now accomplish. It is no longer enough to reopen a preserved deeper relay if the manuscript merely reproduces the same downstream output without addressing the smaller live datum itself. The present note therefore does not return to the ambient package as if the older completion-core theorem still fixed the live burden by itself. Instead it targets the current burden directly: it proves that the smaller trace core \(\TraceCore\) is itself a canonical and necessary reduction of completion-balance ambient dynamics.

This is the cleanest honest way to spend the preferred option left open by the current routing note. The theorem does not enlarge the benchmark package, does not introduce a second operative metric, and does not upgrade adoption to completed first-principles derivation. Its narrower benchmark-facing role is only to show that the present trace-core burden is already forced by a more primitive explicit substrate package. Once that is done, the burden moves below the trace core itself.

\paragraph{Framework and claim boundary.}
\TheoryProgram{} denotes the broader \Aether{} / \Aflow{} framework built on the claims that the \Aether{} is the underlying four-dimensional substrate of reality, the \Aflow{} is its intrinsic ordered motion, observed three-dimensional space is the local experiential slice of that deeper substrate, S-time is the experienced order of change arising from matter, light, and the \Aflow{}, and observed expansion is the three-dimensional appearance of deeper four-dimensional ordered motion. Within that framework, \TheoryName{} denotes the exact-closure theory in which the effective gravitational dynamics are adopted to be exactly Einsteinian with universal matter coupling. In the active sequence, \emph{adoption} means use of that established relativistic dynamics without claiming substrate derivation, while \emph{derivation} is reserved for a first-principles recovery from explicit substrate variables.

The present note stays strictly inside that boundary. It proves only that the current live trace-core burden is a canonical output of a deeper completion-balance ambient dynamics package. It does not claim that the ambient package has already been derived from ultimate substrate microphysics, and it does not license stronger first-principles promotion language.

\section{Fixed Primitive Continuation Data}

\begin{definition}[Recorded primitive continuation data]
\label{def:fixed}
Fix the following continuation data from the active primitive-reservoir observer-reduction line.
\begin{enumerate}
    \item For each sector label \(\sigma\in\{\Car,\Cur,\Hom\}\), a primitive forcing shell
    \[
    \ForSpace{\sigma},
    \]
    a visible reduction map
    \[
    \Reduction{\sigma}:\ForSpace{\sigma}\to X_{\sigma},
    \]
    and shell-level current and equilibrium carriers
    \[
    \CurrentCarrier{\sigma}:\ForSpace{\sigma}\to Z_{\sigma}^{\mathrm{cur}},
    \qquad
    \EqCarrier{\sigma}:\ForSpace{\sigma}\to Z_{\sigma}^{\mathrm{eq}}.
    \]
    \item The product shell
    \[
    \PrimShell
    \equiv
    \ForSpace{\Car}\times \ForSpace{\Cur}\times \ForSpace{\Hom}.
    \]
    \item The recorded metric and ordinary-matter outputs
    \[
    \MetricOut:\PrimShell\to\{\text{observer metrics}\},
    \qquad
    \MatterOut:\PrimShell\times\Psi\to\{\text{observer stress tensors}\},
    \]
    satisfying
    \begin{equation}
    \MetricOut(\mathbf w)=\CompletedMetric,
    \qquad
    \MatterOut(\mathbf w,\psi)
    =
    T_{\mu\nu}^{\mathrm{matter}}[\CompletedMetric,\psi]
    \label{eq:fixedoutputs}
    \end{equation}
    for every \(\mathbf w\in\PrimShell\) and every matter configuration \(\psi\in\Psi\).
\end{enumerate}
\end{definition}

\begin{remark}
Definition \ref{def:fixed} is not reopened here. The primitive shell data, the one operative metric, and the ordinary matter route remain fixed continuation data. The present note only settles the current benchmark-facing burden below the trace-core frontier.
\end{remark}

\section{Current Live Burden: Minimal Completion-State Shell-Trace Core}

\begin{definition}[Completion-state shell-trace core]
\label{def:trace}
Fix \(\sigma\in\{\Car,\Cur,\Hom\}\). A \emph{minimal completion-state shell-trace core with hidden-sector gap} for sector \(\sigma\) consists of the following data.
\begin{enumerate}
    \item A finite-dimensional Euclidean completion-state space
    \[
    \CoreState{\sigma},
    \]
    the observer-data space
    \[
    \DataSpace{\sigma}
    \equiv
    \VisibleAffine{\sigma}\times Z_{\sigma}^{\mathrm{cur}}\times Z_{\sigma}^{\mathrm{eq}},
    \]
    the shell-response space
    \[
    \AmbientTarget{\sigma},
    \]
    and a hidden space
    \[
    \GapSpace{\sigma}.
    \]
    \item A compact convex completion-state body
    \[
    \TraceBody{\sigma}\subset\CoreState{\sigma}
    \]
    with nonempty interior in \(\CoreState{\sigma}\).
    \item A compact convex hidden body
    \[
    \GapBody{\sigma}\subset\GapSpace{\sigma}
    \]
    with \(0\in\operatorname{int}\GapBody{\sigma}\), together with a positive self-adjoint hidden-sector operator
    \[
    \HiddenOperator{\sigma}:\GapSpace{\sigma}\to\GapSpace{\sigma}
    \]
    satisfying
    \begin{equation}
    \ev{\eta,\HiddenOperator{\sigma}\eta}
    \ge
    \kappa_{\sigma}^{\mathrm{hid}}\|\eta\|^{2}
    \qquad
    \text{for all }\eta\in\GapSpace{\sigma}.
    \label{eq:tracegap}
    \end{equation}
    \item Affine observer-data and shell-response maps
    \[
    \TraceDataMap{\sigma}:\CoreState{\sigma}\to\DataSpace{\sigma},
    \qquad
    \TraceShellMap{\sigma}:\CoreState{\sigma}\to\AmbientTarget{\sigma}.
    \]
    \item A compact symmetry group
    \[
    \TraceSym{\sigma}\subset
    O(\CoreState{\sigma})\times
    O(\DataSpace{\sigma})\times
    O(\AmbientTarget{\sigma})\times
    O(\GapSpace{\sigma})
    \]
    and an orthogonal involution
    \[
    \TraceTime{\sigma}\in
    O(\CoreState{\sigma})\times
    O(\DataSpace{\sigma})\times
    O(\AmbientTarget{\sigma})\times
    O(\GapSpace{\sigma})
    \]
    preserving \(\TraceBody{\sigma}\), the exact zero-response trace slice
    \[
    \TraceBody{\sigma}\cap
    \bigl(\TraceShellMap{\sigma}\bigr)^{-1}(0),
    \]
    the hidden body \(\GapBody{\sigma}\), and \(\HiddenOperator{\sigma}\), and satisfying
    \[
    \TraceDataMap{\sigma}(g_{\mathrm{st}}m)
    =
    g_{\mathrm{dat}}\TraceDataMap{\sigma}(m),
    \qquad
    \TraceShellMap{\sigma}(g_{\mathrm{st}}m)
    =
    g_{\mathrm{sh}}\TraceShellMap{\sigma}(m)
    \]
    for every \(m\in\CoreState{\sigma}\) and every
    \[
    (g_{\mathrm{st}},g_{\mathrm{dat}},g_{\mathrm{sh}},g_{\mathrm{hid}})
    \in\TraceSym{\sigma},
    \]
    with the analogous intertwining equalities for the components of \(\TraceTime{\sigma}\).
    \item A smooth embedding
    \[
    \jmath_{\sigma}:X_{\sigma}\hookrightarrow\VisibleAffine{\sigma}
    \]
    and a smooth exact zero-response trace realization
    \[
    \TraceEmbed{\sigma}:\ForSpace{\sigma}\hookrightarrow\CoreState{\sigma}
    \]
    whose image is exactly
    \[
    \TraceBody{\sigma}\cap
    \bigl(\TraceShellMap{\sigma}\bigr)^{-1}(0),
    \]
    and such that
    \begin{equation}
    \TraceDataMap{\sigma}\bigl(\TraceEmbed{\sigma}(w)\bigr)
    =
    \bigl(
    \jmath_{\sigma}(\Reduction{\sigma}(w)),
    \CurrentCarrier{\sigma}(w),
    \EqCarrier{\sigma}(w)
    \bigr)
    \label{eq:tracecompat}
    \end{equation}
    for every \(w\in\ForSpace{\sigma}\).
    \item Fixed-data solvability on the trace body:
    \begin{equation}
    \TraceDataMap{\sigma}\bigl(\TraceBody{\sigma}\bigr)
    =
    \TraceDataMap{\sigma}\Bigl(
    \TraceBody{\sigma}\cap
    \bigl(\TraceShellMap{\sigma}\bigr)^{-1}(0)
    \Bigr).
    \label{eq:tracesolvability}
    \end{equation}
    \item Trace-kernel coercivity: there exists \(\kappa_{\sigma}^{\mathrm{tr}}>0\) such that for every
    \[
    \delta m\in\CoreState{\sigma}
    \]
    satisfying
    \[
    D\TraceDataMap{\sigma}\,\delta m=0,
    \]
    one has
    \begin{equation}
    \bigl\|
    D\TraceShellMap{\sigma}\,\delta m
    \bigr\|^{2}
    \ge
    \kappa_{\sigma}^{\mathrm{tr}}\|\delta m\|^{2}.
    \label{eq:tracecoercive}
    \end{equation}
\end{enumerate}
\end{definition}

\begin{remark}
Definition \ref{def:trace} is the current live internal burden on the active line. The present note asks whether those trace-level data are themselves forced by a still more primitive explicit substrate package, without reintroducing the larger completion-core graph presentation as the stopping point.
\end{remark}

\section{Deeper Explicit Package: Completion-Balance Ambient Dynamics}

\begin{definition}[Completion-balance ambient dynamics with inert completion kernel]
\label{def:ambient}
Fix \(\sigma\in\{\Car,\Cur,\Hom\}\). A \emph{completion-balance ambient dynamics package} for sector \(\sigma\) consists of the following data.
\begin{enumerate}
    \item A finite-dimensional Euclidean ambient state space
    \[
    \BalSpace{\sigma},
    \]
    the observer-data space
    \[
    \DataSpace{\sigma}
    \equiv
    \VisibleAffine{\sigma}\times Z_{\sigma}^{\mathrm{cur}}\times Z_{\sigma}^{\mathrm{eq}},
    \]
    an ambient shell-response space
    \[
    \AmbientTarget{\sigma},
    \]
    and a compact convex ambient admissible body
    \[
    \BalBody{\sigma}\subset\BalSpace{\sigma}
    \]
    with nonempty interior.
    \item An affine observer-data readout
    \[
    \BalData{\sigma}:\BalSpace{\sigma}\to\DataSpace{\sigma}
    \]
    and an affine shell-response readout
    \[
    \BalShell{\sigma}:\BalSpace{\sigma}\to\AmbientTarget{\sigma}.
    \]
    \item An affine zero-hidden slice
    \[
    \ZeroHiddenSlice{\sigma}\subset\BalSpace{\sigma}
    \]
    with translation space
    \[
    \ZeroHiddenVec{\sigma}\equiv T_{z}\ZeroHiddenSlice{\sigma}
    \qquad
    (z\in\ZeroHiddenSlice{\sigma}),
    \]
    and hidden complement
    \[
    \GapSpace{\sigma}
    \equiv
    {\ZeroHiddenVec{\sigma}}^{\perp}\subset\BalSpace{\sigma}.
    \]
    \item A positive self-adjoint ambient balance operator
    \[
    \BalOp{\sigma}:\BalSpace{\sigma}\to\BalSpace{\sigma}
    \]
    whose kernel is exactly the zero-hidden translation space,
    \[
    \ker \BalOp{\sigma}=\ZeroHiddenVec{\sigma},
    \]
    and whose restriction to the hidden complement obeys
    \begin{equation}
    \ev{\eta,\BalOp{\sigma}\eta}
    \ge
    \kappa_{\sigma}^{\mathrm{hid}}\,
    \|\eta\|^{2}
    \qquad
    \text{for all }\eta\in\GapSpace{\sigma}
    \label{eq:ambientgap}
    \end{equation}
    for some \(\kappa_{\sigma}^{\mathrm{hid}}>0\).
    \item The hidden projection body
    \[
    \GapBody{\sigma}
    \equiv
    \left\{
    \eta\in\GapSpace{\sigma}:
    z+\eta\in\BalBody{\sigma}
    \text{ for some }z\in \BalBody{\sigma}\cap\ZeroHiddenSlice{\sigma}
    \right\}
    \]
    is compact and convex, satisfies
    \[
    0\in\operatorname{int}\GapBody{\sigma},
    \]
    and the ambient admissible body splits exactly as
    \begin{equation}
    \BalBody{\sigma}
    =
    \left\{
    z+\eta:
    z\in\BalBody{\sigma}\cap\ZeroHiddenSlice{\sigma},
    \quad
    \eta\in\GapBody{\sigma}
    \right\}.
    \label{eq:bodysplitting}
    \end{equation}
    \item A compact symmetry group
    \[
    \BalSym{\sigma}\subset
    O(\BalSpace{\sigma})\times
    O(\DataSpace{\sigma})\times
    O(\AmbientTarget{\sigma})
    \]
    and an orthogonal involution
    \[
    \BalTime{\sigma}\in
    O(\BalSpace{\sigma})\times
    O(\DataSpace{\sigma})\times
    O(\AmbientTarget{\sigma})
    \]
    preserving \(\BalBody{\sigma}\), \(\ZeroHiddenSlice{\sigma}\), \(\BalData{\sigma}\), \(\BalShell{\sigma}\), and \(\BalOp{\sigma}\).
    \item A smooth embedding
    \[
    \jmath_{\sigma}:X_{\sigma}\hookrightarrow\VisibleAffine{\sigma}
    \]
    and a smooth zero-response shell realization
    \[
    \BalEmbed{\sigma}:\ForSpace{\sigma}\hookrightarrow\BalSpace{\sigma}
    \]
    whose image is exactly
    \[
    \BalBody{\sigma}\cap
    \ZeroHiddenSlice{\sigma}\cap
    \bigl(\BalShell{\sigma}\bigr)^{-1}(0),
    \]
    and such that
    \begin{equation}
    \BalData{\sigma}\bigl(\BalEmbed{\sigma}(w)\bigr)
    =
    \bigl(
    \jmath_{\sigma}(\Reduction{\sigma}(w)),
    \CurrentCarrier{\sigma}(w),
    \EqCarrier{\sigma}(w)
    \bigr)
    \label{eq:ambientcompat}
    \end{equation}
    for every \(w\in\ForSpace{\sigma}\).
    \item Fixed-data shell solvability on the zero-hidden slice:
    \begin{equation}
    \BalData{\sigma}\bigl(
    \BalBody{\sigma}\cap\ZeroHiddenSlice{\sigma}
    \bigr)
    =
    \BalData{\sigma}\Bigl(
    \BalBody{\sigma}\cap
    \ZeroHiddenSlice{\sigma}\cap
    \bigl(\BalShell{\sigma}\bigr)^{-1}(0)
    \Bigr).
    \label{eq:ambientsolvability}
    \end{equation}
    \item The inert completion kernel
    \[
    \Kernel{\sigma}
    \equiv
    \ker\bigl(D\BalData{\sigma}\big|_{\ZeroHiddenVec{\sigma}}\bigr)
    \cap
    \ker\bigl(D\BalShell{\sigma}\big|_{\ZeroHiddenVec{\sigma}}\bigr)
    \]
    and the orthogonal completion-state projector
    \[
    \StateComp{\sigma}:\ZeroHiddenVec{\sigma}\to{\Kernel{\sigma}}^{\perp}\cap\ZeroHiddenVec{\sigma}.
    \]
    \item Completion-balance coercivity on data-invisible zero-hidden directions: there exists \(\kappa_{\sigma}^{\mathrm{co}}>0\) such that for every
    \[
    w\in\ForSpace{\sigma},
    \qquad
    \delta z\in T_{\BalEmbed{\sigma}(w)}\ZeroHiddenSlice{\sigma}=\ZeroHiddenVec{\sigma}
    \]
    satisfying
    \[
    D\BalData{\sigma}\,\delta z=0,
    \]
    one has
    \begin{equation}
    \|D\BalShell{\sigma}\,\delta z\|^{2}
    \ge
    \kappa_{\sigma}^{\mathrm{co}}\,
    \|\StateComp{\sigma}\delta z\|^{2}.
    \label{eq:ambientcoercive}
    \end{equation}
\end{enumerate}
\end{definition}

\begin{remark}
Definition \ref{def:ambient} is deeper than the trace core in exactly the way needed now. It begins with an ambient admissible body, an invariant zero-hidden slice, and an inert kernel on that slice. The current live trace datum must then be recovered as the faithful zero-hidden quotient together with the descended ambient readouts.
\end{remark}

\section{Canonical Trace Core Induced by Ambient Dynamics}

\begin{definition}[Canonical trace core induced by ambient dynamics]
\label{def:canonical}
Assume Definition \ref{def:ambient}. Choose once and for all a reference shell point
\[
w_{\sigma}^{\ast}\in\ForSpace{\sigma},
\qquad
\ReferencePoint{\sigma}\equiv \BalEmbed{\sigma}(w_{\sigma}^{\ast}).
\]
Define the canonical trace core as follows.
\begin{enumerate}
    \item The completion-state space is
    \[
    \CoreState{\sigma}
    \equiv
    {\Kernel{\sigma}}^{\perp}\cap\ZeroHiddenVec{\sigma}.
    \]
    \item The zero-hidden trace coordinate is the affine map
    \begin{equation}
    \TraceCoord{\sigma}(z)
    \equiv
    \StateComp{\sigma}\bigl(z-\ReferencePoint{\sigma}\bigr),
    \qquad
    z\in\ZeroHiddenSlice{\sigma}.
    \label{eq:tracecoord}
    \end{equation}
    \item The completion-state trace body is
    \begin{equation}
    \TraceBody{\sigma}
    \equiv
    \left\{
    \TraceCoord{\sigma}(z):
    z\in\BalBody{\sigma}\cap\ZeroHiddenSlice{\sigma}
    \right\}.
    \label{eq:ambienttracebody}
    \end{equation}
    \item The affine observer-data and shell-response trace maps are
    \begin{equation}
    \TraceDataMap{\sigma}(m)
    \equiv
    \BalData{\sigma}\bigl(\ReferencePoint{\sigma}+m\bigr),
    \qquad
    \TraceShellMap{\sigma}(m)
    \equiv
    \BalShell{\sigma}\bigl(\ReferencePoint{\sigma}+m\bigr)
    \label{eq:ambienttracemaps}
    \end{equation}
    for \(m\in\CoreState{\sigma}\).
    \item The hidden operator is the restriction
    \begin{equation}
    \HiddenOperator{\sigma}
    \equiv
    \BalOp{\sigma}\big|_{\GapSpace{\sigma}}.
    \label{eq:ambienttraceop}
    \end{equation}
    \item The induced symmetry group and time reversal are obtained by restricting the ambient \(\BalSpace{\sigma}\)-action to \(\CoreState{\sigma}\oplus\GapSpace{\sigma}\) and keeping the recorded data and shell actions:
    \[
    \TraceSym{\sigma}
    \equiv
    \left\{
    (g_{\mathrm{tr}},g_{\mathrm{dat}},g_{\mathrm{sh}},g_{\mathrm{hid}}):
    (g_{\mathrm{amb}},g_{\mathrm{dat}},g_{\mathrm{sh}})
    \in\BalSym{\sigma}
    \right\},
    \]
    where \(g_{\mathrm{tr}}\) and \(g_{\mathrm{hid}}\) are the restrictions of \(g_{\mathrm{amb}}\) to \(\CoreState{\sigma}\) and \(\GapSpace{\sigma}\), and
    \[
    \TraceTime{\sigma}
    \equiv
    \bigl(
    \BalTime{\sigma}\big|_{\CoreState{\sigma}},
    \BalTime{\sigma}\big|_{\DataSpace{\sigma}},
    \BalTime{\sigma}\big|_{\AmbientTarget{\sigma}},
    \BalTime{\sigma}\big|_{\GapSpace{\sigma}}
    \bigr).
    \]
    \item The exact zero-response trace realization is
    \begin{equation}
    \TraceEmbed{\sigma}(w)
    \equiv
    \TraceCoord{\sigma}\bigl(\BalEmbed{\sigma}(w)\bigr).
    \label{eq:ambienttraceembed}
    \end{equation}
\end{enumerate}
\end{definition}

\begin{remark}
The chosen reference shell point \(\ReferencePoint{\sigma}\) fixes only the affine origin of the completion-state coordinate. Changing it translates the trace coordinate on the same faithful quotient and does not alter any benchmark-relevant constituent.
\end{remark}

\begin{proposition}[Ambient dynamics induce a trace core]
\label{prop:induce}
Assume Definitions \ref{def:ambient} and \ref{def:canonical}. Then the canonical data of Definition \ref{def:canonical} satisfy every requirement of Definition \ref{def:trace}.
\end{proposition}

\begin{proof}
Because \(\Kernel{\sigma}\subset\ZeroHiddenVec{\sigma}\) and \(\CoreState{\sigma}={\Kernel{\sigma}}^{\perp}\cap\ZeroHiddenVec{\sigma}\), every \(z\in\ZeroHiddenSlice{\sigma}\) can be written uniquely as
\[
z=\ReferencePoint{\sigma}+m+k
\qquad
\text{with }
m\in\CoreState{\sigma},
\quad
k\in\Kernel{\sigma}.
\]
Since \(\BalData{\sigma}\) and \(\BalShell{\sigma}\) are affine and \(\Kernel{\sigma}\) lies in the kernels of their linear parts on \(\ZeroHiddenVec{\sigma}\), the definitions \eqref{eq:ambienttracemaps} are well defined and affine.

The set \(\BalBody{\sigma}\cap\ZeroHiddenSlice{\sigma}\) is compact and convex, so its image \(\TraceBody{\sigma}\) under the affine map \(\TraceCoord{\sigma}\) is compact and convex. Because \(\BalBody{\sigma}\) has nonempty interior and splits exactly as in \eqref{eq:bodysplitting} with \(0\in\operatorname{int}\GapBody{\sigma}\), the zero-hidden slice carries nonempty relative interior in the ambient body. Therefore \(\TraceBody{\sigma}\) has nonempty interior in \(\CoreState{\sigma}\).

Equation \eqref{eq:ambienttraceop} together with \eqref{eq:ambientgap} gives
\[
\ev{\eta,\HiddenOperator{\sigma}\eta}
=
\ev{\eta,\BalOp{\sigma}\eta}
\ge
\kappa_{\sigma}^{\mathrm{hid}}\|\eta\|^{2}
\qquad
(\eta\in\GapSpace{\sigma}),
\]
which is \eqref{eq:tracegap}. The ambient symmetries preserve \(\ZeroHiddenSlice{\sigma}\), \(\ZeroHiddenVec{\sigma}\), \(\GapSpace{\sigma}\), \(\Kernel{\sigma}\), \(\BalBody{\sigma}\), \(\BalData{\sigma}\), \(\BalShell{\sigma}\), and \(\BalOp{\sigma}\), so they induce the trace symmetry group and time reversal, preserving \(\TraceBody{\sigma}\), the exact zero-response trace slice, \(\GapBody{\sigma}\), and \(\HiddenOperator{\sigma}\).

For the exact zero-response realization, \(\TraceEmbed{\sigma}(\ForSpace{\sigma})\) is contained in
\[
\TraceBody{\sigma}\cap
\bigl(\TraceShellMap{\sigma}\bigr)^{-1}(0)
\]
because \(\BalEmbed{\sigma}(\ForSpace{\sigma})\subset
\BalBody{\sigma}\cap\ZeroHiddenSlice{\sigma}\cap(\BalShell{\sigma})^{-1}(0)\). Conversely, if
\[
m\in\TraceBody{\sigma}\cap
\bigl(\TraceShellMap{\sigma}\bigr)^{-1}(0),
\]
then there exists
\[
z\in\BalBody{\sigma}\cap\ZeroHiddenSlice{\sigma}
\]
with \(m=\TraceCoord{\sigma}(z)\). Since \(\TraceShellMap{\sigma}(m)=\BalShell{\sigma}(z)=0\), we have
\[
z\in
\BalBody{\sigma}\cap
\ZeroHiddenSlice{\sigma}\cap
\bigl(\BalShell{\sigma}\bigr)^{-1}(0).
\]
Because \(\BalEmbed{\sigma}\) is a diffeomorphism onto that exact shell-zero locus, there is a unique \(w\in\ForSpace{\sigma}\) with \(z=\BalEmbed{\sigma}(w)\), and therefore \(m=\TraceEmbed{\sigma}(w)\). Equation \eqref{eq:tracecompat} follows directly from \eqref{eq:ambientcompat}.

For fixed-data solvability, let \(u\in\TraceDataMap{\sigma}(\TraceBody{\sigma})\). Then \(u=\BalData{\sigma}(z)\) for some
\[
z\in\BalBody{\sigma}\cap\ZeroHiddenSlice{\sigma}.
\]
By \eqref{eq:ambientsolvability}, there exists
\[
z_{0}\in
\BalBody{\sigma}\cap
\ZeroHiddenSlice{\sigma}\cap
\bigl(\BalShell{\sigma}\bigr)^{-1}(0)
\]
with \(\BalData{\sigma}(z_{0})=u\). Setting \(m_{0}\equiv\TraceCoord{\sigma}(z_{0})\) gives
\[
m_{0}\in
\TraceBody{\sigma}\cap
\bigl(\TraceShellMap{\sigma}\bigr)^{-1}(0)
\]
and \(\TraceDataMap{\sigma}(m_{0})=u\), proving \eqref{eq:tracesolvability}.

Finally, let \(w\in\ForSpace{\sigma}\) and let \(\delta m\in\CoreState{\sigma}\) satisfy
\[
D\TraceDataMap{\sigma}\,\delta m=0.
\]
Viewing \(\delta m\) as a vector of \(\ZeroHiddenVec{\sigma}\), the ambient coercivity inequality \eqref{eq:ambientcoercive} applies because \(\StateComp{\sigma}\delta m=\delta m\). Hence
\[
\|D\TraceShellMap{\sigma}\,\delta m\|^{2}
=
\|D\BalShell{\sigma}\,\delta m\|^{2}
\ge
\kappa_{\sigma}^{\mathrm{co}}\|\delta m\|^{2}.
\]
This is \eqref{eq:tracecoercive} with \(\kappa_{\sigma}^{\mathrm{tr}}=\kappa_{\sigma}^{\mathrm{co}}\).
\end{proof}

\section{Forced Trace Constituents and Uniqueness}

\begin{proposition}[All trace-core constituents are forced by ambient dynamics]
\label{prop:forced}
Assume Definition \ref{def:ambient}. Then the following constituents are fixed by the ambient package itself.
\begin{enumerate}
    \item The completion-state space is the faithful quotient of the zero-hidden translation space by the inert kernel, represented canonically by
    \[
    \CoreState{\sigma}={\Kernel{\sigma}}^{\perp}\cap\ZeroHiddenVec{\sigma}.
    \]
    \item The completion-state trace body is necessarily the image \eqref{eq:ambienttracebody} of the zero-hidden ambient body.
    \item The affine observer-data and shell-response trace maps are necessarily the descended affine readouts \eqref{eq:ambienttracemaps} of the ambient package on the faithful zero-hidden quotient.
    \item The hidden factor is necessarily the orthogonal complement
    \[
    \GapSpace{\sigma}={\ZeroHiddenVec{\sigma}}^{\perp},
    \]
    the hidden body is necessarily \(\GapBody{\sigma}\), and the positive hidden operator is necessarily the restriction \eqref{eq:ambienttraceop} of the ambient balance operator.
    \item The symmetry / time-reversal structure, the exact zero-response trace realization, the fixed-data solvability property, and the trace-kernel coercivity mechanism are all induced by the same ambient package and are not extra trace-level freedom.
\end{enumerate}
\end{proposition}

\begin{proof}
Items (1) through (4) are exactly the canonical constructions of Definition \ref{def:canonical}. Once the zero-hidden slice, its inert kernel, the ambient body, and the ambient readouts are fixed, the only faithful completion-state coordinates are those that quotient out the inert kernel, the only trace body compatible with the same ambient body is \eqref{eq:ambienttracebody}, and the only compatible trace maps are the descended ambient readouts \eqref{eq:ambienttracemaps}. The hidden package is forced by the orthogonal splitting along \(\ZeroHiddenVec{\sigma}\oplus\GapSpace{\sigma}\) together with the balance-operator kernel statement \(\ker\BalOp{\sigma}=\ZeroHiddenVec{\sigma}\).

Item (5) follows because the ambient symmetry group and time reversal preserve the entire package from which those trace constituents are induced, and Proposition \ref{prop:induce} already showed how the exact zero-response realization, fixed-data solvability, and coercivity descend directly to the trace level.
\end{proof}

\section{Downstream Shell Coordinates from the Ambient-Induced Trace Core}

\begin{definition}[Canonical shell coordinates from the trace core]
\label{def:shell}
Assume Definition \ref{def:trace}. Define the shell coordinates induced by the trace core as follows.
\begin{enumerate}
    \item Let
    \[
    \TraceShellLin{\sigma}:\CoreState{\sigma}\to\AmbientTarget{\sigma}
    \]
    denote the linear part of the affine map \(\TraceShellMap{\sigma}\). Define
    \[
    \ShellKer{\sigma}\equiv\ker\TraceShellLin{\sigma},
    \qquad
    \ShellNormal{\sigma}\equiv\operatorname{im}\TraceShellLin{\sigma}.
    \]
    Let
    \[
    \DirProj{\sigma}:\CoreState{\sigma}\to\ShellKer{\sigma}
    \]
    be the orthogonal projection onto \(\ShellKer{\sigma}\), and equip \(\ShellNormal{\sigma}\) with the Euclidean structure transported from \((\ShellKer{\sigma})^{\perp}\) by
    \[
    \TraceShellLin{\sigma}\big|_{(\ShellKer{\sigma})^{\perp}}
    :
    (\ShellKer{\sigma})^{\perp}\xrightarrow{\cong}\ShellNormal{\sigma}.
    \]
    \item Choose the shell offset \(\ShellOffset{\sigma}\in\ShellNormal{\sigma}\) by the affine decomposition
    \begin{equation}
    \TraceShellMap{\sigma}(m)
    =
    \TraceShellLin{\sigma}(m)-\ShellOffset{\sigma}.
    \label{eq:traceshelldecomp}
    \end{equation}
    \item The shell-coordinate space is
    \[
    \CoordSpace{\sigma}\equiv
    \ShellKer{\sigma}\oplus\ShellNormal{\sigma},
    \]
    with product Euclidean structure, and the shell-coordinate isomorphism is
    \begin{equation}
    \CoordIso{\sigma}(m)
    \equiv
    \bigl(
    \DirProj{\sigma}m,\,
    \TraceShellLin{\sigma}(m)
    \bigr).
    \label{eq:tracecoordiso}
    \end{equation}
    \item The shell-coordinate body and shell-data readout are
    \begin{equation}
    \CoordBody{\sigma}
    \equiv
    \CoordIso{\sigma}\bigl(\TraceBody{\sigma}\bigr),
    \qquad
    \AffineData{\sigma}
    \equiv
    \TraceDataMap{\sigma}\circ\CoordIso{\sigma}^{-1}.
    \label{eq:tracecoordbody}
    \end{equation}
    \item The exact shell affine plane is
    \begin{equation}
    \ShellAffine{\sigma}
    \equiv
    \left\{
    (v,e)\in\CoordSpace{\sigma}:
    e=\ShellOffset{\sigma}
    \right\},
    \label{eq:traceshellplane}
    \end{equation}
    and the exact shell realization is
    \begin{equation}
    \ShellEmbed{\sigma}(w)
    \equiv
    \CoordIso{\sigma}\bigl(\TraceEmbed{\sigma}(w)\bigr).
    \label{eq:traceshellembed}
    \end{equation}
    \item The induced symmetry group and time reversal are
    \[
    \SymGroup{\sigma}
    \equiv
    \left\{
    \bigl(
    (g_{\mathrm{st}}|_{\ShellKer{\sigma}},\,g_{\mathrm{sh}}|_{\ShellNormal{\sigma}}),
    g_{\mathrm{hid}}
    \bigr):
    (g_{\mathrm{st}},g_{\mathrm{dat}},g_{\mathrm{sh}},g_{\mathrm{hid}})
    \in\TraceSym{\sigma}
    \right\},
    \]
    \[
    \TimeRev{\sigma}
    \equiv
    \bigl(
    (\TraceTime{\sigma}|_{\ShellKer{\sigma}},
    \TraceTime{\sigma}|_{\ShellNormal{\sigma}}),
    \TraceTime{\sigma}|_{\GapSpace{\sigma}}
    \bigr).
    \]
\end{enumerate}
\end{definition}

\begin{proposition}[The current exact-shell affine nucleus is already carried by the ambient-induced trace core]
\label{prop:nucleus}
Assume Definitions \ref{def:trace} and \ref{def:shell}. Then the shell-coordinate body \(\CoordBody{\sigma}\), affine shell-data readout \(\AffineData{\sigma}\), exact shell affine plane \(\ShellAffine{\sigma}\), exact shell realization \(\ShellEmbed{\sigma}\), hidden body \(\GapBody{\sigma}\), hidden operator \(\HiddenOperator{\sigma}\), induced symmetry / time-reversal structure, fixed-data shell solvability, and affine-kernel coercive control are all determined by the trace core itself. In particular, the shell-coordinate construction yielding the current exact-shell affine nucleus \(\AffineNucleus\) factors through the ambient-induced trace core.
\end{proposition}

\begin{proof}
Because \(\TraceShellLin{\sigma}\big|_{(\ShellKer{\sigma})^{\perp}}\) is a linear isomorphism onto \(\ShellNormal{\sigma}\), the map \(\CoordIso{\sigma}\) is a linear isomorphism; with the chosen Euclidean structures it is an isometry. Hence \(\CoordBody{\sigma}\) is compact and convex with nonempty interior whenever \(\TraceBody{\sigma}\) is.

If \(w\in\ForSpace{\sigma}\), then \(\TraceEmbed{\sigma}(w)\in\TraceBody{\sigma}\) and \(\TraceShellMap{\sigma}(\TraceEmbed{\sigma}(w))=0\). By \eqref{eq:traceshelldecomp},
\[
\TraceShellLin{\sigma}(\TraceEmbed{\sigma}(w))
=
\ShellOffset{\sigma},
\]
so \(\ShellEmbed{\sigma}(w)\in\CoordBody{\sigma}\cap\ShellAffine{\sigma}\). Conversely, if \(q=\CoordIso{\sigma}(m)\in\CoordBody{\sigma}\cap\ShellAffine{\sigma}\), then \(m\in\TraceBody{\sigma}\) and the second coordinate condition forces \(\TraceShellLin{\sigma}(m)=\ShellOffset{\sigma}\), hence \(\TraceShellMap{\sigma}(m)=0\). Therefore \(m=\TraceEmbed{\sigma}(w)\) for a unique \(w\in\ForSpace{\sigma}\), and so \(q=\ShellEmbed{\sigma}(w)\). Thus \(\ShellEmbed{\sigma}\) is a diffeomorphism onto the full exact shell slice of \(\CoordBody{\sigma}\).

Equation \eqref{eq:tracecompat} immediately gives
\[
\AffineData{\sigma}\bigl(\ShellEmbed{\sigma}(w)\bigr)
=
\TraceDataMap{\sigma}\bigl(\TraceEmbed{\sigma}(w)\bigr)
=
\bigl(
\jmath_{\sigma}(\Reduction{\sigma}(w)),
\CurrentCarrier{\sigma}(w),
\EqCarrier{\sigma}(w)
\bigr).
\]
The hidden body, hidden operator, and induced symmetry / time-reversal structure are already part of the trace core.

For shell solvability, let \(u\in\AffineData{\sigma}(\CoordBody{\sigma})\). Then \(u=\TraceDataMap{\sigma}(m)\) for some \(m\in\TraceBody{\sigma}\). By \eqref{eq:tracesolvability}, there exists
\[
m_{0}\in
\TraceBody{\sigma}\cap
\bigl(\TraceShellMap{\sigma}\bigr)^{-1}(0)
\]
with \(\TraceDataMap{\sigma}(m_{0})=u\). Hence
\[
\CoordIso{\sigma}(m_{0})\in\CoordBody{\sigma}\cap\ShellAffine{\sigma}
\]
and \(\AffineData{\sigma}(\CoordIso{\sigma}(m_{0}))=u\).

Finally, let
\[
\delta q\in\CoordSpace{\sigma}
\]
satisfy \(D\AffineData{\sigma}\,\delta q=0\), and let \(\delta m\equiv\CoordIso{\sigma}^{-1}(\delta q)\). Since \(\CoordIso{\sigma}\) is an isometry,
\[
\|\delta q\|=\|\delta m\|.
\]
Moreover,
\[
D\AffineData{\sigma}\,\delta q
=
D\TraceDataMap{\sigma}\,\delta m,
\]
so the trace-kernel condition gives
\[
\bigl\|
D\TraceShellMap{\sigma}\,\delta m
\bigr\|^{2}
\ge
\kappa_{\sigma}^{\mathrm{tr}}\|\delta m\|^{2}.
\]
But the second component of \(\delta q\) is exactly \(D\TraceShellMap{\sigma}\,\delta m\). Therefore the shell-normal part of \(\delta q\) obeys the same coercive lower bound, and adjoining any hidden variation \(\delta n\in\GapSpace{\sigma}\) gives the same affine-kernel coercive control as in the current nucleus theorem. Thus every downstream shell-coordinate constituent relevant to \(\AffineNucleus\) is already carried by the trace core.
\end{proof}

\section{Main Theorems}

\begin{theorem}[Minimal completion-state shell-trace core necessity / uniqueness from ambient dynamics]
\label{thm:necessity}
Assume Definition \ref{def:ambient}. Then:
\begin{enumerate}
    \item the canonical construction of Definition \ref{def:canonical} yields a minimal completion-state shell-trace core in the sense of Definition \ref{def:trace};
    \item every remaining trace-level object singled out by the current trace-core frontier is forced by the ambient package: the completion-state quotient, the compact convex trace body, the affine observer-data and shell-response trace maps, the hidden factor and positive hidden operator, the symmetry / time-reversal structure, the exact zero-response trace realization, the fixed-data solvability property, and the trace-kernel coercivity mechanism;
    \item any other trace core compatible with the same ambient package differs from the canonical one only by the harmless choice of completion-state origin and by an orthogonal coordinate identification on \(\CoreState{\sigma}\);
    \item consequently the shell-coordinate construction yielding the current exact-shell affine nucleus \(\AffineNucleus\) is already carried by the same ambient-induced trace core and does not require the larger completion-core graph presentation as an independently unexplained stopping point.
\end{enumerate}
\end{theorem}

\begin{proof}
Item (1) is Proposition \ref{prop:induce}. Item (2) is Proposition \ref{prop:forced}. For item (3), changing the reference shell point \(\ReferencePoint{\sigma}\) only translates the completion-state coordinate, while changing the faithful orthogonal coordinates on \({\Kernel{\sigma}}^{\perp}\cap\ZeroHiddenVec{\sigma}\) acts only by an orthogonal identification of the same Euclidean completion-state space. No other compatible freedom remains once the ambient zero-hidden slice, inert kernel, ambient body, and ambient readouts are fixed. Item (4) is Proposition \ref{prop:nucleus} together with items (1) and (2).
\end{proof}

\begin{theorem}[The derived trace core preserves the same one-metric benchmark package]
\label{thm:outputs}
Assume Definition \ref{def:ambient}. Then the benchmark output statements of Definition \ref{def:fixed} remain unchanged:
\[
\MetricOut(\mathbf w)=\CompletedMetric,
\qquad
\MatterOut(\mathbf w,\psi)
=
T_{\mu\nu}^{\mathrm{matter}}[\CompletedMetric,\psi]
\]
for every \(\mathbf w\in\PrimShell\) and every matter configuration \(\psi\in\Psi\). Consequently the deeper ambient package, the induced trace core, and the downstream shell-coordinate construction introduce neither a second operative metric nor an enlarged low-energy matter law.
\end{theorem}

\begin{proof}
The present note changes only the upstream origin of the current live trace datum. It does not alter the primitive shell \(\PrimShell\), the recorded metric map \(\MetricOut\), or the recorded matter map \(\MatterOut\). Those remain the fixed continuation data of Definition \ref{def:fixed}, so \eqref{eq:fixedoutputs} continues to hold exactly. The induced trace core and the downstream shell-coordinate construction are built to preserve the same exact shell realization of that unchanged primitive shell, not to replace it with a different observer-output law.
\end{proof}

\begin{theorem}[The remaining burden moves below the trace core]
\label{thm:main}
Assume the fixed primitive continuation data of Definition \ref{def:fixed} and, for each sector \(\sigma\in\{\Car,\Cur,\Hom\}\), a completion-balance ambient dynamics package in the sense of Definition \ref{def:ambient}. Then:
\begin{enumerate}
    \item the minimal completion-state shell-trace core of Definition \ref{def:trace} is canonically induced by the ambient package;
    \item the completion-state quotient, compact convex trace body, affine observer-data and shell-response trace maps, hidden factor and positive hidden operator, symmetry / time-reversal structure, exact zero-response trace realization, fixed-data solvability property, and trace-kernel coercivity mechanism are all forced by that same ambient package;
    \item the shell-coordinate construction yielding the current exact-shell affine nucleus \(\AffineNucleus\) already factors through that ambient-induced trace core;
    \item the one operative metric remains exactly \(\CompletedMetric\), the ordinary matter route remains exactly the standard one through that metric, there is no second operative metric, and the low-energy matter law is not enlarged;
    \item consequently the benchmark-facing burden after the trace-core collapse theorem is discharged in the stronger form now available on the active line: the trace core is no longer an independently unexplained stopping point, but a necessary reduction of a deeper completion-balance ambient dynamics package;
    \item the next honest benchmark-facing attempt should therefore go one level below the current ambient-dynamics frontier: first try to derive or uniquely force that ambient package from still more primitive explicit substrate structure in a way that actually changes benchmark-facing status; if that cannot be done cleanly, the admissible fallback is to prove a genuinely smaller theorem-level collapse, sharper necessity result, or sharper uniqueness result strictly inside that ambient package while preserving the same benchmark one-metric package and claim boundary.
\end{enumerate}
\end{theorem}

\begin{proof}
Item (1) is Theorem \ref{thm:necessity}(1). Item (2) is Theorem \ref{thm:necessity}(2). Item (3) is Theorem \ref{thm:necessity}(4). Item (4) is Theorem \ref{thm:outputs}. Items (5) and (6) are the combined routing consequence.
\end{proof}

\begin{corollary}[What remains open after this theorem]
\label{cor:next}
After Theorem \ref{thm:main}, the remaining honest burden is no longer to justify the trace core itself on the active primitive-observer-reduction line. The live burden is now one level below the current ambient-dynamics frontier. The next honest attempt should therefore first try to derive or uniquely force the completion-balance ambient dynamics package itself from still more primitive explicit substrate structure in a way that actually changes benchmark-facing status while preserving the same benchmark one-metric package and claim boundary. If that cannot be done cleanly, the admissible fallback is to prove a genuinely smaller theorem-level collapse, sharper necessity result, or sharper uniqueness result strictly inside that ambient package.
\end{corollary}

\begin{proof}
Theorem \ref{thm:main} converts the trace core from the current live benchmark-facing burden into a canonical reduction of deeper ambient dynamics. What remains open is why that ambient package itself should hold, rather than becoming the new disciplined stopping point. Since another trace-core restatement, completion-core restatement, exact-shell affine nucleus restatement, or another ambient relay that preserves the same ambient package without shrinking or forcing it further would leave that burden unchanged, the next honest attempt must target the ambient package itself unless a sharper internal collapse, necessity result, or uniqueness result is the only clean benchmark-facing gain available.
\end{proof}

\section{Conclusion}

The positive result of this note is a direct derivation of the current live internal datum from deeper explicit substrate structure. The previous trace-core collapse theorem had already shown that the full completion-core graph presentation carries no extra benchmark-facing freedom beyond the smaller completion-state shell-trace core. The present note then takes the preferred honest next move left open there: it shows that this smaller trace core is itself canonically induced by completion-balance ambient dynamics.

That gain directly addresses the precise constituents that had become the live burden: the completion-state quotient of the zero-hidden slice, the compact convex trace body, the affine observer-data and shell-response trace maps, the hidden body with positive operator, the induced reversible structure, the exact zero-response trace realization, the fixed-data solvability property, and the trace-kernel coercive control. Within the recorded ambient-dynamics class those are no longer optional inserted ingredients. They are forced outputs, unique up to the harmless affine choice of completion-state origin and orthogonal completion-state coordinates.

The benchmark boundary remains fixed throughout. The same primitive shell remains the shell carrying the observer outputs, so the one operative metric is still exactly \(\CompletedMetric\), the ordinary matter route is unchanged, there is no second operative metric, and the low-energy matter law is not enlarged. This is therefore a disciplined upstream derivational gain, not a final first-principles completion claim. The note does not yet derive the ambient completion-balance package from ultimate substrate microphysics. The live burden is now one level below the current ambient-dynamics frontier: first try to derive or uniquely force the ambient package itself from still more primitive explicit substrate structure in a way that actually changes benchmark-facing status. If that cannot be done cleanly, the fallback is a genuinely smaller theorem-level collapse, sharper necessity result, or sharper uniqueness result strictly inside that ambient package. Another trace-core restatement, completion-core restatement, exact-shell affine nucleus restatement, or another ambient relay that preserves the same ambient package without shrinking or forcing it further would not meet that burden, and no stronger first-principles promotion language is licensed yet.

\end{document}
